\indent \textbf{ಉದಾ :}\par

\begin{kbvderiv}
\vspace{4pt plus 3pt minus 1pt}
\begin{center}\begin{adjustbox}{max width=\textwidth}
{\renewcommand{\arraystretch}{1.22}\setlength{\tabcolsep}{4pt}
\begin{tabular}{l c l}
ತಲೆಯ + ಮುಂದು & {\color{kbvgrey}$=$} & \kbvresult{ಮುಂದಲೆ} \\
ಕೈಯ + ಅಡಿ & {\color{kbvgrey}$=$} & \kbvresult{ಅಂಗೈ} \\
ಕಾಲು + ಅಡಿ & {\color{kbvgrey}$=$} & \kbvresult{ಅಂಗಾಲು} \\
ಕಣ್ಣ + ಕಡೆ & {\color{kbvgrey}$=$} & \kbvresult{ಕಡೆಗಣ್ಣು} \\
ಮೈಯ + ಒಳಗು & {\color{kbvgrey}$=$} & \kbvresult{ಒಳಮೈ} \\
ಇರುಳಿನ + ನಡು & {\color{kbvgrey}$=$} & \kbvresult{ನಟ್ಟಿರುಳು} \\
ತುಟಿಯ + ಕೆಳಗು & {\color{kbvgrey}$=$} & \kbvresult{ಕೆಳದುಟಿ} \\
\end{tabular}}
\end{adjustbox}\end{center}
\vspace{4pt plus 3pt minus 1pt}
\end{kbvderiv}

\vspace{4pt plus 4pt minus 1pt}\noindent{\kbvsecfont\fontsize{11pt}{14.5pt}\selectfont\color{kbvaccentdark} (೫)\hspace{0.45em}ದ್ವಂದ್ವ ಸಮಾಸ :}\hspace{0.5em}“ಮಳೆ ಬಂದು ಕೆರೆಕಟ್ಟೆ ತುಂಬಿತು.”— ಈ ವಾಕ್ಯದಲ್ಲಿ ದಪ್ಪಕ್ಷರಗಳಲ್ಲಿರುವ “ಕೆರೆಕಟ್ಟೆ” ಎಂಬ ಪದವನ್ನು ಕೆರೆಯೂ ಕಟ್ಟೆಯೂ ಎಂದು ಬಿಡಿಸಿ ಬರೆಯಬಹುದು. ಕೆರೆಯೂ ಕಟ್ಟೆಯೂ – ಕೆರೆಕಟ್ಟೆ ಇಲ್ಲಿ ಕೆರೆಯೂ ಮತ್ತು ಕಟ್ಟೆಯೂ ಒಂದಕ್ಕೊಂದು ಪೂರಕವಾಗಿರುವ ಅಥವಾ ಸಹಯೋಗ ಹೊಂದಿರುವ ಎರಡು ಪದಗಳು. ಇವುಗಳಲ್ಲಿ ಯಾವು ದೊಂದು ಪದದ ಅರ್ಥವು ಗೌಣವಲ್ಲ, ಎರಡು ಪದಗಳ ಅರ್ಥವೂ ಪ್ರಧಾನವೇ ಸರಿ. ಇಂತಹ ಪದಗಳು ಸೇರಿ ಆಗುವ ಸಮಾಸಕ್ಕೆ ದ್ವಂದ್ವ ಸಮಾಸವೆನ್ನುವರು.\par

\begin{kbvrule}
\noindent{\kbvsecfont ವ್ಯಾಖ್ಯೆ :}\hspace{0.5em}- ಸಹಯೋಗವಿರುವ ಎರಡು ಅಥವಾ ಹೆಚ್ಚು ನಾಮ ಪದಗಳು ಒಂದಕ್ಕೊಂದು ಸೇರಿ ಎಲ್ಲ ಪದಗಳ ಅರ್ಥವೂ ಪ್ರಧಾನವಾಗಿಯುಳ್ಳ ಸಮಾಸಕ್ಕೆ ದ್ವಂದ್ವ ಸಮಾಸವೆಂದು ಹೆಸರು.
\end{kbvrule}

\begin{kbvderiv}
\vspace{4pt plus 3pt minus 1pt}
\begin{center}\begin{adjustbox}{max width=\textwidth}
{\renewcommand{\arraystretch}{1.22}\setlength{\tabcolsep}{4pt}
\begin{tabular}{l c l}
ಉದಾ : ತಾಯಿಯೂ-ತಂದೆಯೂ & {\color{kbvgrey}$=$} & \kbvresult{ತಾಯಿತಂದೆ} \\
ಹಣ್ಣನ್ನೂ-ಕಾಯಿಯನ್ನೂ & {\color{kbvgrey}$=$} & \kbvresult{ಹಣ್ಣುಕಾಯಿ} \\
ಹಸುವನ್ನೂ-ಕರುವನ್ನೂ & {\color{kbvgrey}$=$} & \kbvresult{ಹಸುಕರು} \\
\end{tabular}}
\end{adjustbox}\end{center}
\vspace{4pt plus 3pt minus 1pt}
\end{kbvderiv}
